\clearpage
\SourceLeaf{019}\pdfbookmark[0]{Liber primus.}{liber-primus}
\begin{center}
{\Large LIBER PRIMUS.}\par\bigskip
{\large RELATIONES GENERALES INTER QUANTITATES PER QUAS CORPORUM COELESTIUM MOTUS CIRCA SOLEM DEFINIUNTUR.}\par\bigskip
\pdfbookmark[1]{Sectio prima.}{sectio-prima}
{\large SECTIO PRIMA.}\par\medskip
\emph{Relationes ad locum simplicem in orbita spectantes.}\par\bigskip
\pdfbookmark[2]{1.}{article-001}1.
\end{center}
Corporum coelestium motus in hoc opere eatenus tantum considerabimus, quatenus a Solis vi attractiva gubernantur. Excluduntur itaque ab instituto nostro omnes planetae secundarii, excluduntur perturbationes, quas primarii in se invicem exercent, excluditur omnis motus rotatorius. Corpora mota ipsa ut puncta mathematica spectamus, motusque omnes ad normam legum sequentium fieri supponimus, quae igitur pro basi omnium disquisitionum in hoc opere sunt habendae.

I. Motus cuiusvis corporis coelestis perpetuo fit in eodem plano, in quo simul centrum Solis est situm.

II. Traiectoria a corpore descripta est sectio conica focum in centro Solis habens.

III. Motus in ista traiectoria fit ita, ut areae spatiorum in diversis
\SourceLeaf{020}temporum intervallis circa Solem descriptorum hisce intervallis ipsis sint proportionales. Temporibus igitur et spatiis per numeros expressis, spatium quodvis per tempus intra quod describitur divisum quotientem invariabilem suppeditat.

IV. Pro corporibus diversis circa Solem se moventibus horum quotientium quadrata sunt in ratione composita parametrorum orbitis respondentium, atque aggregatorum massae Solis cum massis corporum motorum.

Designando itaque per $2p$ parametrum orbitae, in qua corpus incedit, per $\mu$ quantitatem materiae huius corporis (posita massa Solis $=1$), per $\frac{1}{2}g$ aream quam tempore $t$ circa Solem describit, erit $\frac{g}{t\BareRoot{p}\cdot\BareRoot{(1+\mu)}}$ numerus pro omnibus corporibus coelestibus constans. Quum igitur nihil intersit, quonam corpore ad valorem huius numeri determinandum utamur, e motu terrae eum depromemus, cuius distantiam mediam a Sole pro unitate distantiarum adoptabimus: unitas temporum semper nobis erit dies medius solaris. Denotando porro per $\pi$ rationem circumferentiae circuli ad diametrum, area ellipsis integrae a terra descriptae manifesto erit $\pi\BareRoot{p}$, quae igitur poni debet $=\frac{1}{2}g$, si pro $t$ accipitur annus sideralis, quo pacto constans nostra fit $=\frac{2\pi}{t\BareRoot{(1+\mu)}}$. Ad valorem numericum huius constantis, in sequentibus per $k$ denotandae, explorandum, statuemus, secundum novissimam determinationem, annum sideralem sive $t=365{,}2563835$, massam terrae sive $\mu=\frac{1}{354710}=0{,}000\,0028\,192$, unde prodit

\begin{center}
\begin{tabular}{@{}ll@{}}
\makebox[15em][l]{$\log 2\pi$\hfill .\,.\,.\,.\,.\,.\,.\,.\,.\,.\,.\,.\,.\,.} & $0{,}798\,1798\,684$\\[.3em]
\makebox[15em][l]{Compl. $\log t$\hfill .\,.\,.\,.\,.\,.\,.\,.\,.\,.} & $7{,}437\,4021\,852$\\[.3em]
\makebox[15em][l]{Compl. $\log\BareRoot{(1+\mu)}$\hfill .\,.\,.\,.\,.} & $9{,}999\,9993\,878$\\[.3em]
\hline
\makebox[15em][l]{$\log k$\hfill .\,.\,.\,.\,.\,.\,.\,.\,.\,.\,.\,.\,.\,.\,.} & $8{,}235\,5814\,414$\\[.3em]
\multicolumn{1}{c}{$k=$} & $0{,}017\,2020\,9895$
\end{tabular}
\end{center}
